\uuid{BVeh}
\titre{Intégration par parties}
\auteur{hanani ; mzari}
\organisation{exo7}
\licence{CC BY-NC-SA 4.0}
\datecreate{2020-01-20}
\niveau{L1}
\module{Analyse}
\chapitre{Calcul d'intégrales}
\sousChapitre{Primitives diverses}
\difficulte{1}
\theme{Primitives, intégration par parties}
\exo7qcmid{lille-2.primitives.009}

\contenu{
\begin{question}
\qclassification{127.11}
\qtags{motcle=intégration par parties}
On note par $F$ une primitive de $f(x)=x\mathrm{e}^x$ sur $\Rr$. Parmi les affirmations suivantes, cocher celles qui sont vraies :
\begin{answers}
\bad{$\displaystyle F(x)=x\ln x+k$, où $k\in \Rr$.}
\bad{$\displaystyle F(x)=\frac{x^2}{2}\mathrm{e}^x+k$, où $k\in \Rr$.}
\good{$\displaystyle F(x)=x\mathrm{e}^x-\int \mathrm{e}^x\mathrm{d}x$.}
\good{$\displaystyle F(x)=(x-1)\mathrm{e}^x+k$, où $k\in \Rr$.}
\end{answers}
\begin{explanations}
Une intégration par parties avec $u=x$ et $v=\mathrm{e}^x$ donne
$$F(x)=x\mathrm{e}^x-\int \mathrm{e}^x\mathrm{d}x=x\mathrm{e}^x-\mathrm{e}^x+k,\mbox{ où }k\in \Rr.$$
\end{explanations}
\end{question}
}
